
\documentclass[11pt,a4paper]{article}
\usepackage[utf8]{inputenc}
\usepackage[T1]{fontenc}
\usepackage[spanish]{babel}
\usepackage{geometry}
\usepackage{graphicx}
\usepackage{array}
\usepackage{tabularx}
\usepackage{makecell}
\usepackage{hyperref}
\usepackage{xcolor}
\usepackage{fancyhdr}
\usepackage{titlesec}
\usepackage{helvet}
\usepackage{ragged2e}
\usepackage{enumitem}
\usepackage{setspace}
\usepackage{amsmath}
\usepackage{booktabs}
\usepackage{float}
\usepackage{listings}
\usepackage{hyperref}
\usepackage{xurl}
\usepackage{microtype}
\usepackage{float}
\emergencystretch=2em
\renewcommand{\familydefault}{\sfdefault}
\geometry{a4paper,left=3cm,right=3cm,top=2.35cm,bottom=2.35cm,headheight=1.9cm,headsep=0.45cm}
\setlength{\parindent}{0pt}
\setlength{\parskip}{0.3em}
\setlength{\tabcolsep}{5pt}
\renewcommand{\arraystretch}{1.13}
\pagestyle{fancy}
\fancyhf{}
\fancyhead[L]{\IfFileExists{reporte_tecnico_overleaf_logo.png}{\includegraphics[height=1.35cm]{reporte_tecnico_overleaf_logo.png}}{\small Universidad Nacional de Loja}}
\fancyhead[R]{\small FEIRNNR -- Carrera de Computación}
\fancyfoot[C]{\thepage}
\renewcommand{\headrulewidth}{0.5pt}
\renewcommand{\footrulewidth}{0pt}
\titleformat{\section}{\bfseries\fontsize{15}{17}\selectfont}{\thesection.}{0.8em}{}
\titlespacing*{\section}{0pt}{1.0ex}{0.3ex}
\newcolumntype{L}[1]{>{\RaggedRight\arraybackslash}p{#1}}
\newcolumntype{Y}{>{\RaggedRight\arraybackslash}X}
\lstset{basicstyle=\ttfamily\footnotesize,breaklines=true,frame=single,showstringspaces=false,columns=fullflexible}
\hypersetup{hidelinks}

\begin{document}
\thispagestyle{fancy}
\begin{center}
{\bfseries\fontsize{21}{24}\selectfont Reporte Técnico de Actividades\\ Práctico-Experimentales Nro. 004\par}
\vspace{0.25cm}
{\large Análisis comparativo de herramientas de simulación mediante un modelo común}
\end{center}

\section{Datos de Identificación del Estudiante y la Práctica}
\begin{tabularx}{\textwidth}{|L{0.43\textwidth}|Y|}
\hline
\textbf{Nombre del estudiante(s)} & [Miguel Alejandro Armas ordoñez] \\ \hline
\textbf{Asignatura} & Simulación \\ \hline
\textbf{Ciclo} & 5 \\ \hline
\textbf{Unidad} & U1: Introducción a la Simulación \\ \hline
\textbf{Resultado de aprendizaje} & Aplica análisis estadísticos de los datos de la simulación para mejorar la calidad de los resultados, bajo los principios de solidaridad, transparencia, responsabilidad y honestidad. \\ \hline
\textbf{Práctica Nro.} & 004 \\ \hline
\textbf{Título de la Práctica} & Análisis comparativo de herramientas de simulación mediante un modelo común \\ \hline
\textbf{Nombre del Docente} & José O. Guamán Q. \\ \hline
\textbf{Fecha} & 22 y 23 de octubre de 2026 \\ \hline
\textbf{Horario} & Jueves 10h30--11h30; viernes 07h30--09h30 \\ \hline
\textbf{Lugar} & Aula \\ \hline
\textbf{Tiempo planificado en el Sílabo} & 3 horas \\ \hline
\end{tabularx}

\section{Objetivo(s) de la Práctica}


\section{Materiales, Reactivos, Equipos y Herramientas}


\section{Procedimiento / Metodología Ejecutada}


\section{Resultados}

\section{Preguntas de Control}

\section{Conclusiones}


\section{Recomendaciones}


\section{Bibliografía / Referencias}

\section{Anexos}

\section{Declaración de uso de Inteligencia Artificial}

Durante el desarrollo de esta práctica se utilizó ChatGPT como herramienta de apoyo para la implementación y revisión del modelo de simulación SIR. Su uso estuvo enfocado principalmente en la generación de gráficas para representar los resultados y en la orientación para desarrollar el código en NetLogo, debido al desconocimiento inicial de su sintaxis, estructura y funcionamiento.

Asimismo, se solicitó apoyo para verificar el código implementado, identificar posibles errores y reconocer aspectos que podían mejorarse, especialmente en la organización modular, la aplicación de las fórmulas matemáticas y la exportación de resultados en archivos CSV.

La inteligencia artificial se utilizó como una herramienta complementaria de consulta y asistencia técnica. Su participación no reemplazó el análisis, la interpretación de resultados ni la toma de decisiones durante el desarrollo de la práctica.

\end{document}
